% Part: history
% Chapter: biographies
% Section: ernst-zermelo

\documentclass[../../../include/open-logic-section]{subfiles}

\begin{document}

\olfileid{his}{bio}{zer}

\olsection{اېرنست زېر مېلو}

\olphoto{zermelo-ernst}{اېرنست زېر مېلو}

اېرنست زېر مېلو د 1871 د جولای پر 27مه
د جرمني په برلين کښې زېږېدلے و۔
هغه پنځۀ خويندې لرلې، خو کورنۍ
ئې د خرابې روغتيا له امله کړېده
او يوازې درې ئې تر بالغ عمره
ژوندۍ پاتې شوې۔ مور او پلار ئې
هم په کم عمرۍ کښې وفات شول؛
کله چې هغه اووۀ‌لس کلن و، هغه
او وروڼه خويندې ئې يتيمان شول۔
زېر مېلو په هنرونو، او په تېره
په شاعرۍ، کښې ژور شوق لرلو۔
هغه په ځيرکتيا، حاضرځوابۍ او
نيوکه‌کوونکي خوي پېژندل کېده۔
د هغه په تر ټولو مشهورو رياضيکي
لاسته راوړنو کښې په 1904 کښې
د ټاکنې د بديهي اصل معرفي کول
او په 1908 کښې د سټونو د
تيورۍ اصل‌بندي کول شامل دي۔

په پوهنتون کښې د زېر مېلو شوقونه
بېلابېل وو۔ هغه د فزيکس،
رياضياتو او فلسفې درسونه واخيستل۔
د هېرمان شوارڅ تر لارښوونې لاندې
ئې په 1894 کښې د برلين په
پوهنتون کښې خپله رساله
\emph{د تغيراتو په حساب کښې
څېړنې} بشپړه کړه۔ په 1897
کښې ئې د ګوټينګن په پوهنتون
کښې د نورو زده کړو پرېکړه وکړه،
او هلته د ډېوېډ هېلبرټ بنسټيز
کار ډېر اغېزمن کړ۔ په 1899
کښې د پروفېسرۍ اهل شو، خو
دنده ئې تر يوولسو کلونو وروسته
ترلاسه نۀ کړه؛ ښايي علت ئې
عجيب چلند او «عصبي تلوار» و۔

زېر مېلو بالاخره په 1910 کښې د زيوريخ
په پوهنتون کښې معاش لرونکې
پروفېسري ترلاسه کړه، خو په
1916 کښې د نري رنځ له امله
تقاعد ته اړ شو۔ له رغېدو وروسته
ئې په 1921 کښې د فرايبورګ په
پوهنتون کښې افتخاري پروفېسري
ترلاسه کړه۔ په دې وخت کښې ئې
د رياضياتو پر بنسټونو کار کولو۔
د تورالف سکولېم او کُرټ ګوډل
آثار ئې په غصه کړل، او په خپلو
مقالو کښې ئې په ښکاره د هغوي
پر لارو نيوکې وکړې۔ په 1935
کښې په جرمني کښې د هټلر د
واک ته رسېدو د مخالفت او د
خپل نامقبوليت له امله د
فرايبورګ له دندې لرې کړے شو۔

د زېر مېلو د ژوند وروستي کلونه په
ګوښه‌والي کښې تېر شول۔ په 1935
کښې له دندې له لرې کېدو وروسته
ئې رياضيات پرېښودل۔ کليوالې
سيمې ته ولاړ او په ساده ډول
ئې ژوند کولو۔ په 1944 کښې ئې
وادۀ وکړ، او د سترګو د نظر د
کمېدو له امله بيخي په مېرمن
متکي شو۔ تر 1951 پورې ئې نظر
په بشپړ ډول لاړو۔ هغه د 1953
د مې پر 21مه د جرمني په
ګينترشتال کښې وفات شو۔

\begin{reading}
د زېر مېلو د بشپړ ژوندليک دپاره
\citet{Ebbinghaus2015} وګورئ۔ د
هغه د 1904 او 1908 بنسټيزې
مقالې په اصلي جرمني کښې لوستل
کېدلے شي
\citep{Zermelo1904,Zermelo1908}۔ د
زېر مېلو راټول آثار، چې د فزيکس
ليکنې ئې هم پکښې دي،
په انګليسي ژباړه کښې په
\citep{Ebbinghaus2010,Ebbinghaus2013}
کښې شته۔
\end{reading}

\end{document}
